% =============================================================================
%  Per-axis explanation of V / A / D for the VADBridge report.
%  Drop-in: requires amsmath, amssymb, booktabs.
% =============================================================================

\subsection{The Three Affective Axes}
\label{sec:vad-axes}

The three dimensions are populated by deliberately disjoint families of
kinematic signals -- \emph{shape} for $V$, \emph{dynamics} for $A$, and
\emph{direction} for $D$ -- so that the resulting representation is
near--orthogonal by construction rather than by post--hoc decorrelation.

% -----------------------------------------------------------------------------
\subsubsection*{V $\cdot$ Valence}

\paragraph{Psychological definition.}
The ``pleasant--unpleasant'' polarity of an affective
state~\cite{russell1980circumplex,mehrabian1996pad}. $V>0$ corresponds to
pleasure / approach; $V<0$ to displeasure / withdrawal.

\paragraph{Bodily correlates (shape signals).}
We select three indicators that describe the \emph{form} of the motion,
deliberately excluding any directional or kinetic information:
\begin{enumerate}\itemsep1pt
  \item \emph{flow} vs. brittleness (Laban Effort\,\textperiodcentered\,Flow),
  \item \emph{expansion} vs. contraction (Camurri contraction
        index~\cite{camurri2003recognizing}),
  \item \emph{spine uprightness} vs. forward slump~\cite{boone2001children}.
\end{enumerate}

\paragraph{Formulas.}
\begin{align}
  \phi   &\;=\; \mathbb{1}[\bar s>s_0]\cdot
          \Bigl(1 - \operatorname{clip}\!\bigl(J/(\bar s+\varepsilon),0,1\bigr)\Bigr),
          \label{eq:V-phi}\\[2pt]
  \kappa &\;=\; \frac{1}{T\,J}\sum_{t=1}^{T}\sum_{j=1}^{J}\bigl\|x^{\text{loc}}_{t,j}\bigr\|_2,
          \label{eq:V-kappa}\\[2pt]
  u      &\;=\; 1 - \frac{1}{T}\sum_{t=1}^{T}\max\!\bigl(0,\,-\sin\text{pitch}_t\bigr).
          \label{eq:V-u}
\end{align}
The motion gate $\mathbb{1}[\bar s>s_0]$ in~\eqref{eq:V-phi} ($s_0=0.02$
rad/frame) prevents a frozen pose from being scored as maximally smooth;
$u$ is asymmetric so that backward lean and upright both saturate while
only forward lean penalises valence.

\paragraph{Fusion.}
\begin{equation}
  V \;=\; 0.40\,\tilde\phi + 0.35\,\tilde\kappa + 0.25\,\tilde u.
  \label{eq:V-fuse}
\end{equation}

\paragraph{Difficulty.}
Karg~\textit{et~al.}~\cite{karg2013body} (\S7.2) identify $V$ as the
\emph{hardest} of the three dimensions to recover from body motion,
because the same kinematic envelope can carry opposite valences depending
on micro--shape (e.g.\ smooth vs. brittle execution of an identical
gesture).

\paragraph{Example.}
The same waving gesture yields $V\!\approx\!+0.6$ when executed smoothly,
expansively, and upright (warm greeting), and $V\!\approx\!-0.5$ when
jittery, contracted, and slouched (anxious wave).

% -----------------------------------------------------------------------------
\subsubsection*{A $\cdot$ Arousal}

\paragraph{Psychological definition.}
The ``activation--rest'' intensity of an affective
state~\cite{russell1980circumplex}. $A>0$: tense / energetic / alert;
$A<0$: calm / relaxed / drowsy.

\paragraph{Bodily correlates (dynamics signals).}
We choose three indicators that span the zeroth, first, and second
temporal derivatives of joint angles:
\begin{enumerate}\itemsep1pt
  \item \emph{mean speed} -- ``the most commonly selected feature''
        across the affective--computing literature~\cite{karg2013body},
  \item \emph{jerk} -- whether the motion starts and stops abruptly,
  \item \emph{acceleration peak} -- a proxy for Laban
        Effort\,\textperiodcentered\,Weight (``force'' moments).
\end{enumerate}

\paragraph{Formulas.}
\begin{align}
  \bar s    &\;=\; \frac{1}{29\,T}\sum_{t=1}^{T}\sum_{j=1}^{29}\bigl|\dot q_{t,j}\bigr|,
              \label{eq:A-speed}\\[2pt]
  J         &\;=\; \frac{1}{29\,(T-3)}\sum_{t=1}^{T-3}\sum_{j=1}^{29}
              \bigl|q_{t+3,j}-3q_{t+2,j}+3q_{t+1,j}-q_{t,j}\bigr|,
              \label{eq:A-jerk}\\[2pt]
  a_{\max}  &\;=\; \max_{t,\,j}\bigl|q_{t+2,j}-2q_{t+1,j}+q_{t,j}\bigr|.
              \label{eq:A-accmax}
\end{align}
Note: $J$ uses $L_1$ (not $L_2$) for outlier robustness, and $a_{\max}$
takes the \emph{peak} (not the mean) so that a single ``moment of effort''
separates a stomp from a smooth walk.

\paragraph{Fusion.}
\begin{equation}
  A \;=\; 0.40\,\tilde{\bar s} + 0.35\,\tilde J + 0.25\,\tilde a_{\max}.
  \label{eq:A-fuse}
\end{equation}

\paragraph{Difficulty.}
$A$ is the \emph{easiest} of the three dimensions
(Karg~\textit{et~al.}~\cite{karg2013body} \S7.2 reports the highest
cross--study agreement). This also explains why our augmentation
framework needs only a single temporal--resample primitive
(\S\ref{sec:vad-aug}, Opt~4) to traverse the entire $A$ axis without
touching any DOF.

\paragraph{Example.}
The same walking action: slow and steady $A\!\approx\!-0.4$ (relaxed),
brisk $A\!\approx\!+0.5$ (hurried), jogging $A\!\approx\!+0.7$
(excited).

% -----------------------------------------------------------------------------
\subsubsection*{D $\cdot$ Dominance}

\paragraph{Psychological definition.}
The sense of control over a situation~\cite{mehrabian1996pad}. $D>0$:
agentic / target--directed / leading; $D<0$: passive / hesitant /
yielding.

\paragraph{A deliberate re--framing.}
The classical literature operationalises $D$ via \emph{static expansion}
cues (bounding--box volume, head height, chest puff -- ``I look big,
therefore I am dominant''). VADBridge deliberately rejects this framing
and defines $D$ as \emph{outward action} towards a target, for three
reasons:
\begin{enumerate}\itemsep1pt
  \item \textbf{Data.} Wallbott's PCA of bodily affect~\cite{wallbott1998bodily}
        loads openness entirely onto $V$. Re--using expansion for $D$ would
        guarantee a $V\!\leftrightarrow\!D$ collapse rather than fight it.
  \item \textbf{Story.} The NMI cross--channel headline (gesture +
        handover) requires $D$ to be computable and semantically
        consistent on \emph{both} channels. ``Approach + reach + direct''
        survives both; static expansion survives neither cleanly.
  \item \textbf{Prior art.} Hall's proxemics~\cite{hall1966hidden},
        Burgoon's forward--lean studies, and Tracy \&
        Robins' pride--posture work~\cite{tracy2004show} all anchor $D$
        in \emph{outward direction}, not in static size.
\end{enumerate}
The price -- absorbed deliberately -- is that all three $D$ indicators
require motion: a frozen pose has $D\!\approx\!0$, which is physically
appropriate (a stationary body cannot ``dominate'' anything).

\paragraph{Bodily correlates (direction signals).}
\begin{enumerate}\itemsep1pt
  \item \emph{reach extension} -- wrists forward of the
        pelvis~\cite{ekman1972hand},
  \item \emph{forward approach} -- character--frame translation along the
        facing axis~\cite{hall1966hidden},
  \item \emph{directness} -- net displacement over path length
        (Laban Effort\,\textperiodcentered\,Space~\cite{laban1947effort}).
\end{enumerate}

\paragraph{Formulas.}
\begin{align}
  r              &\;=\; \frac{1}{T}\sum_{t=1}^{T}\max\!\left(0,\;\tfrac{1}{2}\bigl(x^{\text{loc}}_{t,\text{L-wrist},\text{fwd}}+x^{\text{loc}}_{t,\text{R-wrist},\text{fwd}}\bigr)\right),
                    \label{eq:D-reach}\\[2pt]
  v_{\text{fwd}} &\;=\; \frac{1}{T}\sum_{t=1}^{T}\Delta p^{\text{local}}_{t,\text{fwd}},
                    \label{eq:D-vfwd}\\[2pt]
  \delta         &\;=\; \frac{\bigl\|\sum_{t=1}^{T}\Delta p_t\bigr\|_2}{\sum_{t=1}^{T}\bigl\|\Delta p_t\bigr\|_2+\varepsilon}.
                    \label{eq:D-delta}
\end{align}
$v_{\text{fwd}}$ is signed, so a backward step contributes negative
dominance (withdrawal).

\paragraph{Fusion.}
\begin{equation}
  D \;=\; 0.40\,\tilde v_{\text{fwd}} + 0.40\,\tilde r + 0.20\,\tilde\delta.
  \label{eq:D-fuse}
\end{equation}
The equal $0.40$ weight on $r$ and $v_{\text{fwd}}$ deliberately mirrors
the dual--channel design goal: gesture--type clips obtain their $D$
signal predominantly through reach, while locomotion clips obtain it
through forward translation; each channel has a strong anchor. Path
directness ($\delta$) is a Laban Space--Direct modifier carrying lower
weight.

\paragraph{Difficulty.}
$D$ is the most contested of the three
(Karg~\textit{et~al.}~\cite{karg2013body} \S7.4.2 notes it is often
simplified or omitted). Tying $D$ to the social--handover narrative is
one of the paper--grade contributions of the present work.

\paragraph{Augmentation primitive.}
$D$ is modulated via $\mathcal{P}_{\text{lean}}$ (\S\ref{sec:vad-aug},
Opt~5), which rotates the \emph{root quaternion} about its body--$Y$ axis
rather than perturbing the \texttt{waist\_pitch} DOF: the $D_1$
indicator channel reads $\sin(\text{root pitch})$, while
\texttt{waist\_pitch} feeds $V_3$ (uprightness). Modulating the DOF would
shift $V$ without shifting $D$.

\paragraph{Empirical orthogonality.}
On a $10$k--clip BONES sample the Pearson correlation between $V$ and
$D$ is
\begin{equation}
  r(V, D) \;=\; 0.052 \;\approx\; 0,
\end{equation}
confirming that the disjoint indicator partition (shape for $V$,
direction for $D$) yields the intended independence in real data.

\paragraph{Example.}
In a handover scenario: stepping forward and extending the arm first
yields $D\!\approx\!+0.6$ (assertive); standing still with hands at the
sides and waiting for the partner to reach yields $D\!\approx\!-0.4$
(deferent).

% -----------------------------------------------------------------------------
\paragraph{Summary.}
\begin{table}[h]
  \centering\small
  \begin{tabular}{@{}llllc@{}}
    \toprule
    Axis & Signal family & Indicators & Primary augmentation & Recoverability \\
    \midrule
    $V$ & shape    & $\phi,\kappa,u$              & Opt 1 / 2 / 3 (amp, squat, openness) & hardest \\
    $A$ & dynamics & $\bar s,J,a_{\max}$          & Opt 4 (temporal resample)            & easiest \\
    $D$ & direction & $r,v_{\text{fwd}},\delta$    & Opt 5 (forward lean)                 & most contested \\
    \bottomrule
  \end{tabular}
  \caption{The three families of kinematic signals are physically
           disjoint (shape vs. dynamics vs. direction), which makes the
           nine indicators mutually non--competing and the five
           augmentation primitives freely composable.}
  \label{tab:vad-axes-summary}
\end{table}

\endinput
